\section{Special functions, asymptotic forward models, and real thresholds}
\label{sec:special-function-adapters}

The operation \wl{AsymptoticSpecialInverse} exposes explicit adapters for
\wl{"Erfc"}, \wl{"LogGamma"}, \wl{"Gamma"},
\wl{"LambertThreshold"}, and \wl{"QuadraticThreshold"}.
These examples need two different contracts. The tail adapters invert a
finite Poincar\'e forward model and transport its omitted error. The
threshold adapters use exact local equations in a ramified coordinate.
A convergent inverse expansion for a finite model does not make the
original Poincar\'e forward series convergent.

Every adapter permits a real affine change of target,
$y=b+aF(x)$, with a provably real offset and a provably nonzero real scale.
It stores the original function, the forward model or exact local equation,
the transformed target, the selected real branch, and the meaning of the
requested cutoff. The following formulas use $a=1,b=0$ unless stated
otherwise.

\subsection{The complementary-error-function tail}

For $x>0$, define
\begin{equation}
 S(x)=\sqrt\pi\,x e^{x^2}\operatorname{erfc}x,
 \qquad
 S_m(x)=\sum_{k=0}^{m-1}\frac{(-1)^k(1/2)_k}{x^{2k}},
 \qquad A_m=(1/2)_m.
 \label{eq:erfc-normalized-tail}
\end{equation}
The real positive-axis expansion has a remainder bounded by its first
neglected term; see
\href{https://dlmf.nist.gov/7.12.i}{DLMF~\S7.12(i)}.
For inverse-error transport we also require a derivative bound.

\begin{lemma}[Value and derivative bounds for the normalized tail]
\label{lem:erfc-normalized-bounds}
For every integer $m\ge1$ and every $x>0$,
\begin{equation}
 |S(x)-S_m(x)|\le A_m x^{-2m},\qquad
 |S'(x)-S_m'(x)|\le 2m A_m x^{-2m-1}.
 \label{eq:erfc-normalized-bounds}
\end{equation}
\end{lemma}
\begin{proof}
Starting with the integral definition of $\operatorname{erfc}$ in
\href{https://dlmf.nist.gov/7.2.E2}{DLMF~(7.2.2)}, substitute
$s^2=x^2+t$ to obtain
\begin{equation}
 S(x)=\int_0^\infty e^{-t}(1+t/x^2)^{-1/2}\,dt.
 \label{eq:erfc-normalized-integral}
\end{equation}
Put $w=x^{-2}$. The derivatives of $(1+tw)^{-1/2}$ alternate in sign, and
their magnitudes for $w\ge0$ are at most their values at $w=0$.
Taylor's integral remainder, followed by integration in $t$, bounds the
remainder of the first $m$ terms by $(1/2)_m w^m$.
Applying the same argument to the derivative with respect to $w$ bounds
its remainder by $m(1/2)_m w^{m-1}$. For $m=1$, the latter statement is
the direct bound on the first derivative. Finally
$|dw/dx|=2x^{-3}$ gives the second inequality in
\eqref{eq:erfc-normalized-bounds}. The exponential factor in the integral
justifies these finite differentiations and integrations.
\end{proof}

Let $Q_m(w)$ be the Taylor polynomial of degree $m-1$ of
$-\log S_m(w^{-1/2})$ at $w=0$. The adapter constructs the finite phase
\begin{equation}
 \Phi_m(x)=x^2+\log x+Q_m(x^{-2}),
 \qquad v=-\log(\sqrt\pi\,y).
 \label{eq:erfc-finite-phase}
\end{equation}
The exact phase is $\Phi(x)=-\log(\sqrt\pi\operatorname{erfc}x)$.
For fixed $m$, Lemma~\ref{lem:erfc-normalized-bounds} implies
\begin{equation}
 \Phi(x)-\Phi_m(x)=\OO(x^{-2m}),\qquad
 \Phi'(x)-\Phi_m'(x)=\OO(x^{-2m-1}).
 \label{eq:erfc-phase-errors}
\end{equation}
These are value and derivative assertions; the second is not inferred
merely by differentiating an unspecified big-$O$ term.

The implementation records explicit conditional bounds as well. Write
$\delta=A_mx^{-2m}$, $\delta_1=2mA_mx^{-2m-1}$, and
$\widetilde\Phi_m=x^2+\log x-\log S_m$. Wherever $S_m>\delta$,
\begin{align}
 |\Phi-\Phi_m|
 &\le\frac{\delta}{S_m-\delta}
       +|\widetilde\Phi_m-\Phi_m|,
       \label{eq:erfc-explicit-phase-bound}\\
 |\Phi'-\Phi_m'|
 &\le\frac{\delta_1}{S_m-\delta}
 +\frac{|S_m'|\delta}{S_m(S_m-\delta)}
 +|\widetilde\Phi_m'-\Phi_m'|.
 \label{eq:erfc-explicit-phase-derivative-bound}
\end{align}
The terms involving $\widetilde\Phi_m$ are exact elementary expressions.
The positivity condition is retained in the result; it is satisfied
eventually because $S_m\to1$ and $\delta\to0$.

Now invert $\Phi_m(x)=v$ on $x\to+\infty$. Its leading source scale is
$x^2$, so the inverse error transported from \eqref{eq:erfc-phase-errors}
is $\OO(x^{-2m-1})=\OO(v^{-m-1/2})$. The ordinary inverse engine receives
the explicit local input remainder $\{2m,0\}$ in $u=1/x$. A requested
cutoff larger than $m+1/2$ is rejected unless the forward model is refined.
This bound is combined with the truncation error of the model inverse.

For example, with $v=-\log(\sqrt\pi y)$,
\begin{equation}
 \operatorname{erfc}^{-1}y
 =\sqrt v-\frac{\log v}{4\sqrt v}
 +\frac{-\log^2v+4\log v-8}{32v^{3/2}}
 +\OO\!\left(v^{-5/2}(1+|\log v|)^3\right).
 \label{eq:erfc-inverse-first-terms}
\end{equation}
The cubic numerator follows by inserting
$x=\sqrt v+A/\sqrt v+B/v^{3/2}$ into
$x^2+\log x+(2x^2)^{-1}+\OO(x^{-4})=v$.
The expansion has a power cutoff in $1/v$, although $y\to0^+$.

The negative-source tail uses the exact identity
$\operatorname{erfc}(-x)=2-\operatorname{erfc}x$ and changes the
reconstructed source sign. For $b+a\operatorname{erfc}x$ at $-\infty$,
the positive small tail is $2-(y-b)/a$ and the target endpoint is $b+2a$.
These identities handle reflected, shifted and sign-scaled tails without
changing the forward-error proof.

\subsection{A Stirling model with a retained Lambert inverse}

For $x>0$, let
\begin{equation}
 G_m(x)=x(\log x-1)-\frac12\log x+\frac12\log(2\pi)
 +\sum_{k=1}^{m-1}\frac{B_{2k}}{2k(2k-1)x^{2k-1}}.
 \label{eq:stirling-finite-model}
\end{equation}
The positive-axis Stirling and digamma error bounds in
\href{https://dlmf.nist.gov/5.11.ii}{DLMF~\S5.11(ii)} give
\begin{align}
 |\log\Gamma(x)-G_m(x)|&\le A_mx^{-\rho},
 \label{eq:stirling-value-bound}\\
 |\psi(x)-G_m'(x)|&\le\rho A_mx^{-\rho-1},
 \label{eq:stirling-derivative-bound}\\
 \rho&=2m-1,\qquad
 A_m=\frac{|B_{2m}|}{2m(2m-1)}.
 \notag
\end{align}
The finite value and derivative bounds apply for all positive $x$.
The adapter selects the increasing inverse branch $x>2$; on this interval
$\psi(x)\ge\log x-(2x)^{-1}-(12x^2)^{-1}>0$.
This branch choice excludes the separate behavior near the positive
minimum of the gamma function.

Use the exact core
\begin{equation}
 F_0(x)=x(\log x-1),\qquad
 \phi(v)=\frac{v}{W_0(v/e)}.
 \label{eq:stirling-exact-core}
\end{equation}
The difference $G_m-F_0$ is a finite higher-power perturbation in
$u=1/x$. The operation of Section~\ref{sec:core-perturbations} therefore
retains $\phi$ exactly and computes its complete marker corrections.
With $u_0=1/\phi(v)$, the omitted Stirling error transports to
\begin{equation}
 \OO(u_0^{2m-1})
 \label{eq:stirling-inverse-input-floor}
\end{equation}
in the implementation's conservative polynomial--log bound. The actual
derivative of the core provides an additional logarithmic denominator,
but discarding that improvement preserves a valid bound in the stated
remainder class. The final error combines
\eqref{eq:stirling-inverse-input-floor} with the complete marker tail.

For $m=1$, the first marker correction gives
\begin{equation}
 (\log\Gamma)^{-1}(v)
 =\phi(v)+\frac12-\frac{\log(2\pi)}{2\log\phi(v)}
 +\OO(\phi(v)^{-1}).
 \label{eq:stirling-first-corrected-inverse}
\end{equation}
Increasing marker depth while holding $m$ fixed cannot improve the inherited
Stirling accuracy limit. Marker coefficients beyond that limit still
describe the chosen finite model, with that interpretation explicit in
the result. Neither the finite marker expansion nor its analytic
complete-tail theorem establishes convergence of the infinite Stirling
series.

For inversion of $\Gamma(x)$, use the exact target transformation
$v=\log y$. The implementation solves and checks the equation
$\log\Gamma(x)=\log y$; it can therefore accept a target such as
$e^{10000}$ without forming that enormous value in the forward residual.
For a real affine target $b+a\Gamma(x)$, replace $y$ by $(y-b)/a>1$.
The analogous log-gamma adapter uses $v=(y-b)/a>0$ directly.

\subsection{The two real Lambert branches at their common threshold}

The two real inverses of $F(x)=xe^x$ meet at $x=-1$, $y=-1/e$.
The branch description is given in
\href{https://dlmf.nist.gov/4.13}{DLMF~\S4.13}.
Writing $x=-1+h$ gives the exact local equation
\begin{equation}
 eF(-1+h)+1
 =\frac{h^2}{2}+\frac{h^3}{3}+\frac{h^4}{8}
 +\frac{h^5}{30}+\cdots.
 \label{eq:Lambert-threshold-forward}
\end{equation}
The leading power is two. Ordinary analytic reversion in the positive
ramified coordinate
\begin{equation}
 t=\sqrt{2(1+ey)}
 \label{eq:Lambert-threshold-uniformizer}
\end{equation}
therefore yields a convergent local expansion. Its first coefficients are
\begin{equation}
 W_\nu(y)=-1+\sigma t-\frac{t^2}{3}
 +\sigma\frac{11t^3}{72}-\frac{43t^4}{540}+\OO(t^5),
 \label{eq:Lambert-threshold-inverse}
\end{equation}
where $(\nu,\sigma)=(0,1)$ or $(-1,-1)$.
The implementation computes these coefficients from the exact forward
equation, not from a large-argument Lambert asymptotic formula.

The option \wl{"LambertBranch"} selects $0$ or $-1$.
Its source direction must agree: the principal branch approaches $-1$
from above and the lower branch from below. Contradictory choices are
rejected. The recorded real target neighborhood is $-1/e<y<0$,
transformed by the requested affine target map. A target below the real
threshold is rejected by the numerical checker. The cutoff is an
exclusive power cutoff in the positive distance $y+1/e$, so a cutoff
$5/2$ retains the terms through $t^4$.

\subsection{An exact coalescing quadratic branch}

For an exact quadratic model
\begin{equation}
 F(x)=b+a c(x-x_0)^2,\qquad ac\ne0,
 \label{eq:quadratic-threshold-forward}
\end{equation}
the two branches are exactly
\begin{equation}
 x=x_0+\sigma\sqrt{\frac{y-b}{ac}},\qquad
 \sigma\in\{1,-1\},\qquad\frac{y-b}{ac}>0.
 \label{eq:quadratic-threshold-inverse}
\end{equation}
The source direction selects $\sigma$. The remainder is zero because this
is an exact inverse identity. For $F(x)=x^2-\mu$, the branches at target
$y=0$ are $\pm\sqrt\mu$ and coalesce as $\mu\to0^+$. The same coordinate
retains the selected sign throughout this limiting comparison. It does
not choose a different branch because the two roots become close.

\subsection{API, numerical evidence, and switching policy}

\begin{lstlisting}[language=Mathematica]
AsymptoticSpecialInverse["Erfc", {x, Infinity}, {y, 3}]
AsymptoticSpecialInverse["LogGamma", {x, Infinity}, {y, 3}]
AsymptoticSpecialInverse["Gamma", {x, Infinity}, {y, 2}]
AsymptoticSpecialInverse[
  "LambertThreshold", {x, -1}, {y, 5/2},
  "LambertBranch" -> -1]
AsymptoticSpecialInverse[
  "QuadraticThreshold", {x, 0}, {y, 1},
  "TargetOffset" -> -mu, Assumptions -> mu > 0,
  Direction -> "FromBelow"]
\end{lstlisting}

The gamma calls use an integer marker depth. The error-function call uses
a power cutoff in its logarithmic target coordinate. The threshold calls
use local target-power cutoffs. Each result identifies this distinction.
\wl{"ModelTerms"} controls the number of retained forward asymptotic
terms, with a default sufficient for the requested error-function cutoff
or gamma marker depth. A model that cannot support an error-function
cutoff is rejected; a deliberately coarse Stirling model retains its
explicit accuracy floor.

\wl{SpecialInverseNumericalCheck} uses the original special-function
equation or the exact threshold inverse as its independent reference.
Exact target substitutions precede numerical evaluation, preserving
$\log(e^v)=v$ and avoiding unnecessary cancellation in reflected tiny
error-function tails. If a numerical input has already lost the digits
needed by its transformed target, the checker reports insufficient
precision. Its output sets \wl{"Certified"} to \wl{False}: a
high-precision solve is not an interval proof.

The regression families compare successive model orders at common target
values, test the first-neglected value and derivative bounds against the
original special functions, compare local Lambert expansions with exact
\wl{ProductLog} values on both real branches, and use the exact quadratic
formula as an oracle. These are explicit overlap comparisons. They do not
establish a universal automatic crossover threshold between a local
expansion, a tail expansion and another numerical representation.

The current adapters are selected explicitly. Their metadata states this
switching policy. Complex branches, automatic certified crossover
selection, the gamma minimum and its second positive inverse branch,
arbitrary coalescing higher-order critical points, and additional special
functions remain outside these admitted families. The four roadmap
acceptance types are represented here with their own branch and error
contracts; their existence does not make these remaining cases automatic.
